\subsection{All finite cyclic outer groups, with their actual automorphism extensions}

\begin{statement}{D.47}{Proved}{Complete peripheral chief fibres for every cyclic outer group}
Let \(S\) be a finite nonabelian simple group for which
\[
 \epsilon:\operatorname{Aut}(S)\longrightarrow
       \operatorname{Out}(S)=C_M
 \tag{D.47.1}
\]
has cyclic target, where \(M\geq1\) is arbitrary.
Identify \(S\) with its inner automorphisms.
Let \(1\to A=S^\Omega\to E\xrightarrow{\pi}D\to1\)
be an actual finite chief extension, and fix an epimorphism
\(q_E:\widehat F(x_1,x_2)\twoheadrightarrow E\).
Write \(x_3=(x_1x_2)^{-1}\),
\(X_i=q_E(x_i)\), \(q_0=\pi q_E\), and \(d_i=q_0(x_i)\).
For every epimorphism \(q_1:\widehat F\twoheadrightarrow D\),
every specified \(h_i\in D\) with
\[
                    q_1(x_i)=h_id_ih_i^{-1},
 \tag{D.47.2}
\]
and arbitrary lifts \(H_i\in E\) of these \(h_i\), the finite set
\[
 \mathcal E(q_1;H_i)=
 \left\{f:\widehat F\twoheadrightarrow E:
   \pi f=q_1,\quad
   f(x_i)\in\{aH_iX_iH_i^{-1}a^{-1}:a\in A\}
                 \ (1\leq i\leq3)\right\}
 \tag{D.47.3}
\]
has exactly the same positive cardinality as
\(\mathcal E(q_0;1,1,1)\).
This includes the complete ordinary peripheral fibre and every
choice of its prescribed base conjugators.
The actual extension \textup{(D.47.1)} is retained throughout.

Consequently every D.37 automorphism \(\beta_Q\) on every
specified \(P_n\) lifts through every chief extension of this
kind above \(P_n\), with the literal conjugators
\(1,f_Q,h_Q\). Its maximal pro-five quotient is the given
\(P_n\), and the induced map there is \(\beta_Q\).
Any additional prescribed epimorphism to a finite group over
\(D\) with solvable kernel is retained through the actual
fibre product.
\end{statement}

\begin{proof}
We prove the composite-order source maps first, and then give
the exact extension, boundary, subtraction, and generation maps.

\paragraph{Representatives in the actual outer extension.}
Choose an element \(t\in\operatorname{Aut}(S)\) with
\(\epsilon(t)=1\in C_M\). For \(a\in C_M\), choose the
integer representative \(\widetilde a\in\{0,\ldots,M-1\}\),
and put
\[
                  L(a)=t^{\widetilde a},\qquad L(0)=1.
 \tag{D.47.4}
\]
The elements
\(L(a)L(b)L(a+b)^{-1}\) belong to the actual subgroup \(S\).
They remain in all products below. Every change of coordinates
will use a displayed element \(L(a)\) and its actual inverse.
For a subgroup \(H\leq C_M\), write
\(\widetilde S_H=\epsilon^{-1}(H)\).
If \(|H|=m\), its generator \(M/m\) has the actual lift
\(t_H=t^{M/m}\); identify \(H\) with \(C_m\) using that
generator.

\paragraph{The effective phase of a boundary.}
Use the oriented surface presentation of D.45.6, with
\(b\geq2\), and let \(\chi:\Pi\twoheadrightarrow C_m\)
be a surjective character. Write \(j_i=\chi(c_i)\).
An automorphism preserving \(\chi\) and every positive
boundary class has
\(\theta(c_i)=u_ic_iu_i^{-1}\). Its effective phase is
\[
                 [\chi(u_i)]\in C_m/\langle j_i\rangle.
 \tag{D.47.5}
\]
This quotient retains exactly the ambiguity in the based path.
Each boundary is primitive in the free group obtained by
eliminating one other boundary. Its centralizer is its cyclic
subgroup: in a free basis containing that boundary, a reduced
word with a letter outside that cyclic subgroup cannot commute
with the basis letter. Thus two conjugating words differ by
a boundary power, proving \textup{(D.47.5)} is well-defined.
For character-preserving automorphisms the effective phases
add under composition. Under a change of surface coordinates,
transport \(\chi\) with that change; the phases of the change
and its inverse then cancel in these quotients.

The same quotient is exactly what is needed for target
conjugacy sets. For \(a\in\widetilde S_H\) with outer
coordinate \(j\in C_m\), and \(k\in\langle j\rangle\),
choose an integer \(v\) with \(vj=k\) in \(C_m\).
Then \(t_H^ka^{-v}\in S\), and the actual equality
\[
 t_H^k a t_H^{-k}
     =(t_H^ka^{-v})\,a\,(t_H^ka^{-v})^{-1}
 \tag{D.47.6}
\]
shows that every \(S\)-conjugacy invariant subset of that
outer coset is fixed by the phase subgroup \(\langle j\rangle\).
No order relation for \(t_H\) was used.

\paragraph{Every effective phase is realized by source automorphisms.}
For one handle and a boundary \(c\), retain the exact maps
\(P\) and \(T\) on \([a,b]c\):
\[
 \begin{aligned}
 P(a)&=c^{-1}ac,&P(b)&=c^{-1}a^{-1}cab c,&
 P(c)&=c^{-1}a^{-1}cac,\\
 T(a)&=a,&T(b)&=ba,&T(c)&=c.
 \end{aligned}
 \tag{D.47.7}
\]
Both fix \([a,b]c\) literally. D.45.8--9 give the
explicit Nielsen inverse of \(P\), and \(T^{-1}(b)=ba^{-1}\).
They commute. Moreover \(P\) fixes \((ac)^{-1}\), and its
restriction to \(\langle a,c\rangle\) is conjugation by that
word. Hence for all integers \(N,k\), the map
\(\theta=P^N\circ T^k\) is an actual automorphism, with
inverse \(T^{-k}\circ P^{-N}\), and
\[
 \begin{split}
  \chi(\theta(a))&=\chi(a),\qquad
  \chi(\theta(c))=j,\\
  \chi(\theta(b))&=\chi(b)+Nj+k\chi(a),\\
  \theta(c)&=(ac)^{-N}c(ac)^N.
 \end{split}
 \tag{D.47.8}
\]

Fix a prime power \(\ell^e\Vert m\), and let
\(N\in\mathbf Z\) be a CRT idempotent: it is \(1\) modulo
\(\ell^e\) and \(0\) modulo every other primary factor of \(m\).
If \(\ell\nmid j\), the entire \(\ell\)-primary subgroup
of \(C_m\) lies in \(\langle j\rangle\), so this component
already acts trivially by \textup{(D.47.6)}.
Suppose \(\ell\mid j\). Surjectivity of \(\chi\) supplies
a handle or boundary value which is a unit modulo \(\ell\).

First let a handle supply such a value. If its \(a\)-value
is a unit, put \(a_0=\chi(a)\). If only its \(b\)-value is
a unit, use the exact change \(a\mapsto ab,\ b\mapsto b\)
and its inverse; the new \(a_0\) is \(\chi(a)+\chi(b)\),
which is a unit modulo \(\ell\). In these coordinates choose
\[
 k\equiv-j a_0^{-1}\pmod{\ell^e},\qquad
 k\equiv0\pmod{m/\ell^e}.
 \tag{D.47.9}
\]
Equations \textup{(D.47.8)--(D.47.9)} give
\(Nj+ka_0=0\) in \(C_m\), so \(\theta\) preserves
the entire character. Its phase at \(c\) is
\(-N(a_0+j)\). That phase is a unit at \(\ell\) and is
zero at all other primary factors. Every other boundary is
fixed literally. Powers of \(\theta\) therefore supply
every \(\ell\)-primary effective phase at \(c\).

If a boundary \(d\) supplies the unit, put
\(d_0=\chi(d)\). The exact Hurwitz square of D.45.11--12,
raised to the \(N\)-th power, conjugates \(c\) by \((cd)^N\).
It conjugates \(d\) by \((cd)^Nd^{-N}\). These literal words
give phases
\[
                     N(j+d_0)\quad\hbox{at }c,
                         \qquad Nj\quad\hbox{at }d.
 \tag{D.47.10}
\]
The first is a unit in the required primary component and
zero elsewhere. The second belongs to \(\langle d_0\rangle\):
its other primary components vanish, and \(d_0\) is a unit
in the remaining component. Thus its effective phase at
\(d\) is zero. The map preserves \(\chi\), fixes all other
boundaries, and its powers supply the required component at
\(c\).

The positive boundary swaps and handle-block swaps in
D.45.11 and D.45.13 put any selected boundary and supplier
in these positions. Pull back by their explicit inverses.
The comparison following \textup{(D.47.5)} retains the
effective phases in the original presentation. This proves
surjectivity of the effective-phase map onto
\(\prod_i C_m/\langle j_i\rangle\), using finite
compositions of actual surface automorphisms.

It follows that for arbitrary \(S\)-conjugacy invariant
sets \(T_i\subset\widetilde S_H\) in the specified outer
cosets, independently replacing them by
\(t_H^{k_i}T_it_H^{-k_i}\) preserves the number of
epimorphisms \(\widehat\Pi\twoheadrightarrow\widetilde S_H\)
with outer character \(\chi\) and those boundary conditions.
Indeed precomposition by the constructed automorphism gives
the bijection: its conjugating word has image in
\(S t_H^{k_i}\), up to the harmless subgroup in
\textup{(D.47.6)}. Precomposition by the explicit inverse
reverses the bijection and preserves surjectivity.

\paragraph{Matching composite cyclic characters exactly.}
Two oriented pants covers of the same index, with the matched
boundary cycles below, have the same genus. Choose standard
handle presentations and use the positive boundary swaps to
match their boundary classes. Consider two characters to
\(C_m\) which have the same image \(C_m\) and the same
values \(j_1,\ldots,j_b\) on those matched boundaries.
Define
\[
       d=\gcd(m,\widetilde j_1,\ldots,\widetilde j_b),
       \qquad (j_1,\ldots,j_b)=d(\mathbf Z/m)\subset\mathbf Z/m.
 \tag{D.47.11}
\]
In genus zero the characters already agree. In positive
genus, their handle covectors reduce modulo \(d\) to
primitive covectors. Over each local ring
\(\mathbf Z/\ell^a\) dividing \(d\), elementary symplectic
operations carry a primitive covector to the first basis
covector: place a unit entry in the first symplectic pair,
clear the other entries, and scale that pair by
\(\operatorname{diag}(u,u^{-1})\). The elementary
factorization in D.41 realizes the scaling as well.
All these elementary operations are powers of the actual
handle Dehn twists listed there. Choose their parameters
by CRT to be the required ones in the selected primary
factor and zero in the others. Thus a finite product of
integral surface maps matches the handle covectors modulo
\(d\), while fixing every boundary.

The remaining difference in each handle coordinate lies
in the ideal generated by the \(j_i\) in
\textup{(D.47.11)}. The bounding-pair map for a basis curve
\(a\) and a boundary \(c_i\) acts on homology by
\(v\mapsto v+\langle v,a\rangle c_i\), as in D.45.15.
Its powers change the dual handle-character coordinate by
arbitrary multiples of \(j_i\), and fix all the other
handle coordinates and the boundary values. An integral
linear combination expressing the required difference in
the ideal \((j_1,\ldots,j_b)\) therefore corrects that
coordinate exactly. Doing this for each handle matches
the characters in \(C_m\). These are actual based surface
maps, with inverse maps. Their additional boundary phases
are all removed by \textup{(D.47.5)--(D.47.10)}.
For the trivial image the characters agree without this
correction. This proves the required epimorphism-count
comparison for every cyclic image, including every proper
subgroup of \(C_M\).

\paragraph{The given chief extension and each intermediate embedding.}
Let \(D_*\) be a factor stabilizer and use the right action
\(\Omega=D_*\backslash D\). Centerlessness gives the exact
isomorphism
\[
 E\longrightarrow\operatorname{Aut}(A)
             \mathbin{\times_{\operatorname{Out}(A)}}D,
          \qquad e\longmapsto(\operatorname{Ad}_e,\pi(e)).
 \tag{D.47.12}
\]
Injectivity follows from \(Z(A)=1\); for surjectivity,
choose a lift of the prescribed base element and correct
its automorphism by the unique required inner element of
\(A\).
Write the outer action as the actual cocycle
\(b_d(\omega)\in C_M\), with
\(b_{de}(\omega)=b_d(\omega)+b_e(\omega d)\).
Choose \(t_\omega\in D\), \(t_{\omega_0}=1\), and put
\[
 \begin{aligned}
 \epsilon_*(c)&=b_c(\omega_0)\quad(c\in D_*),&
 s_\omega&=b_{t_\omega}(\omega_0),\\
 k(\omega,d)&=t_\omega d t_{\omega d}^{-1},&
 b_d(\omega)&=-s_\omega+\epsilon_*(k(\omega,d))+s_{\omega d}.
 \end{aligned}
 \tag{D.47.13}
\]
The actual gauge
\((z,d)\mapsto(L(s_\omega)z(\omega)L(s_{\omega d})^{-1},d)\)
carries \(E\), and every specified \(X_i,H_i\), to
\[
 E_*=
 \{(z,d)\in\operatorname{Aut}(S)^\Omega\rtimes D:
             \epsilon(z(\omega))=\epsilon_*(k(\omega,d))\},
 \qquad
 (z,d)(w,e)=(z(\omega)w(\omega d),de).
 \tag{D.47.14}
\]
For multiplication, the adjacent factors
\(L(s_{\omega d})^{-1}L(s_{\omega d})\) cancel literally.
This proves the map and its inverse using the actual
representatives \textup{(D.47.4)}.

For each \(D_*\leq C\leq D\), retain every character
\(\epsilon_C:C\to C_M\) extending \(\epsilon_*\).
Choose \(R_\nu\) representing \(C\backslash D\), with
\(R_C=1\). Define
\[
 \begin{aligned}
 k_C(\nu,d)&=R_\nu dR_{\nu d}^{-1},&
 b_C(d,\nu)&=\epsilon_C(k_C(\nu,d)),\\
 p(\omega)&=Ct_\omega,&
 g_\omega&=\epsilon_C(t_\omega R_{p(\omega)}^{-1}).
 \end{aligned}
 \tag{D.47.15}
\]
Let \(E_C\) be the group in \textup{(D.47.14)} with
\(C\backslash D\) and the outer pattern \(b_C\).
There is an actual injection over \(D\)
\[
 j_C:E_C\hookrightarrow E_* ,\qquad
 j_C(z,d)(\omega)=L(g_\omega)z(p(\omega))L(g_{\omega d})^{-1}.
 \tag{D.47.16}
\]
Its outer-coordinate identity is
\(\epsilon_*(k(\omega,d))=
g_\omega+b_C(d,p(\omega))-g_{\omega d}\).
The displayed factors in \(C\) prove this equality;
adjacent inverse factors prove the homomorphism identity.
Reading a coordinate above each \(\nu\) proves injectivity.
Also \(g_{\omega_0}=0\).

\paragraph{The actual source and every boundary root condition.}
Put \(K_j=\ker q_j\) and \(H_{j,C}=q_j^{-1}(C)\).
For words \(T_\omega\) with \(q_j(T_\omega)=t_\omega\)
and \(T_{\omega_0}=1\), twisted Shapiro is the explicit map
\[
 \begin{split}
 \varphi:H_{j,D_*}&\longrightarrow\operatorname{Aut}(S),
                      \qquad \epsilon\varphi=\epsilon_*q_j,\\
 z_\varphi(g)(\omega)
       &=\varphi(T_\omega gT_{\omega q_j(g)}^{-1}),\qquad
 f_\varphi(g)=(z_\varphi(g),q_j(g)).
 \end{split}
 \tag{D.47.17}
\]
Middle words cancel in every product. Conversely evaluate
the cocycle at \(\omega_0\) on \(H_{j,D_*}\).
The normalizing vector
\(v(\omega)=z(T_\omega)(\omega_0)\) belongs to \(S^\Omega\),
has \(v(\omega_0)=1\), and gives
\(z_\varphi(g)(\omega)=
v(\omega)z(g)(\omega)v(\omega q_j(g))^{-1}\).
These formulas are inverse normalizations, so each surface
map has exactly \(|S|^{|\Omega|-1}\) literal gauges.
They preserve the boundary \(A\)-classes and generation.
The analogous construction on \(C\backslash D\), followed
by \(j_C\), restricts at \(\omega_0\) to the very same
surface map on \(H_{j,D_*}\): here \(g_{\omega_0}=0\)
and \(D_*\) fixes that coordinate.

Conjugacy by \(A\) for two elements above the same base
element is equivalent to \(S\)-conjugacy of every
permutation-cycle product. Indeed the successive conjugator
equations determine each next coordinate from the preceding
one; their outer coordinates vanish, and the last equation
is precisely the cycle-product condition.
For a boundary \(\delta\) on \(C\backslash D\), choose
initial coordinate \(\nu\), and coordinates \(\omega\)
over it for all fine cycles \(\gamma\) of relative degree
\(m_\gamma\). The exact enlarged boundary subset
\(T_{j,C,\delta}\) is the intersection
\[
 \begin{gathered}
 a\in\widetilde S_{\operatorname{im}\epsilon_C},\qquad
 \epsilon(a)=\epsilon_C(q_j(c_\delta)),\\
 L(g_\omega)a^{m_\gamma}L(g_\omega)^{-1}
                     \in\mathcal C_{j,\gamma}
                         \quad\hbox{for every }\gamma.
 \end{gathered}
 \tag{D.47.18}
\]
The classes \(\mathcal C_{0,\gamma}\) come from the
actual cycle products of \(X_i\); those for \(j=1\)
come from \(H_iX_iH_i^{-1}\). These subsets retain all
roots and every empty intersection.

The actual boundary bijection between the two covers is
\[
 C d\langle q_1(x_i)\rangle
              \longmapsto C d h_i\langle d_i\rangle.
 \tag{D.47.19}
\]
It preserves lengths and every relative degree.
Take the old initial coordinate corresponding to a new
\(\nu'\) to be \(\nu'h_i\).
There is an actual element \(\mathcal L_i\in E_C\)
above \(h_i\), with components
\(\mathcal L_i(\nu)=L(b_C(h_i,\nu))\).
Its image under \(j_C\) and \(H_i\) differ by an element
of \(A\). Conjugating by \(\mathcal L_i\) therefore gives
the exact equality of complete boundary subsets
\[
 T_{1,C,\delta'}=
 L(b_C(h_i,\nu'))\,T_{0,C,\delta}\,L(b_C(h_i,\nu'))^{-1}.
 \tag{D.47.20}
\]
On a cycle product the conjugator is precisely this single
component of \(\mathcal L_i\). Its inverse proves the
reverse inclusion, including the intersections in
\textup{(D.47.18)}. The argument allows arbitrary given
\(X_i\); membership of \(X_i\) in \(j_C(E_C)\) is not
required for the conjugacy comparison.

Both source characters \(\epsilon_Cq_j|_{H_{j,C}}\)
have the same image \(\operatorname{im}\epsilon_C\),
since \(q_j(H_{j,C})=C\), and their matched boundary
values agree. Their standard surface presentations have
the same genus by \textup{(D.47.19)} and the degree formula.
Apply the exact character-matching maps above, then the
independent effective-phase maps to \textup{(D.47.20)}
and to the paths in the matching isomorphism.
Every such phase lies in \(\operatorname{im}\epsilon_C\).
This proves equality, for \(j=0,1\), of the numbers of
epimorphisms
\[
 H_{j,C}\twoheadrightarrow
         \widetilde S_{\operatorname{im}\epsilon_C}
 \tag{D.47.21}
\]
with their exact outer character and all conditions
\textup{(D.47.18)}.

\paragraph{Subtracting the identical finite-base maps.}
For an epimorphism in \textup{(D.47.21)}, its \(K_j\)-image
is normal in the target and contained in \(S\), hence is
either \(1\) or \(S\). The maps with image \(1\) are
exactly
\[
 \varphi=\Phi q_j|_{H_{j,C}},\qquad
 \Phi:C\twoheadrightarrow\widetilde S_{\operatorname{im}\epsilon_C},
                  \qquad\epsilon\Phi=\epsilon_C.
 \tag{D.47.22}
\]
Use the same actual \(\Phi\) on both sides. Its Shapiro
section over \(D\) is the homomorphism
\[
 s_\Phi:D\longrightarrow E_C,\qquad
 s_\Phi(d)(\nu)=\Phi(R_\nu dR_{\nu d}^{-1}),
               \quad\hbox{with base coordinate }d.
 \tag{D.47.23}
\]
It is independent of \(q_j\). Equations
\textup{(D.47.2)} and \textup{(D.47.23)} show that
\(j_Cs_\Phi(q_1(x_i))\) is the conjugate of
\(j_Cs_\Phi(d_i)\) by \(j_Cs_\Phi(h_i)\).
This conjugator and \(H_i\) have the same base value,
so their ratio belongs to \(A\). Thus the old boundary
condition holds for this \(\Phi\) exactly when the new
one does. This is a bijection of the precise subtracted
sets, and it retains all character-extension choices.
Subtracting their equal finite cardinalities from
\textup{(D.47.21)} gives equal counts of maps with
\(\varphi(K_j)=S\), for every actual intermediate
\(C\) and every extending \(\epsilon_C\).

\paragraph{Normalizer inversion and the entire generating fibre.}
For a surface map on \(H_{j,D_*}\) with \(K_j\)-image \(S\),
let \(N_\varphi=\ker(\varphi|_{K_j})\).
Its normalizer contains \(H_{j,D_*}\) and \(K_j\),
so it equals \(H_{j,C}\) for some \(C\geq D_*\).
The normalizer contains \(H_{j,C}\) precisely when
\(\varphi\) extends to \(H_{j,C}\to\operatorname{Aut}(S)\).
Conjugation on \(K_j/N_\varphi\simeq S\) constructs this
extension. It agrees with \(\varphi\) on \(H_{j,D_*}\)
and is unique: an automorphism centralizing all inner
automorphisms of the centerless group \(S\) fixes every
element of \(S\). Its outer coordinate factors through
one unique extending character \(\epsilon_C:C\to C_M\).
Its image is the full actual preimage
\(\widetilde S_{\operatorname{im}\epsilon_C}\), because
its \(K_j\)-image is \(S\). Conversely such an extension
normalizes \(N_\varphi\).

Let \(M_{j,C}\) be the sum, over every extending
\(\epsilon_C\), of the retained counts after
\textup{(D.47.22)}. The literal restriction comparison
for \(j_C\) proves that this counts exactly those
original surface maps whose kernel normalizer contains
\(H_{j,C}\). Uniqueness makes the summands disjoint.
The preceding paragraphs give \(M_{0,C}=M_{1,C}\).
Finite subgroup-poset inversion therefore gives
\[
 e_j=\sum_{D_*\leq C\leq D}\mu(D_*,C)M_{j,C},\qquad
 e_0=e_1,\qquad
 |\mathcal E(q_j;H_i)|=|S|^{|\Omega|-1}e_j,
 \tag{D.47.24}
\]
where on the old side the conjugators are \(1\).
To verify the last equality, minimal normalizer means
that the conjugate coordinate kernels on \(K_j\) are
pairwise distinct. The \(K_j\)-image is a subdirect
product of the simple factors. A proper such product
identifies two factors and their coordinate kernels:
apply the two-factor Goursat lemma inductively, using
that every normal subgroup of a product of nonabelian
simple groups is a product of factors. Thus minimal
normalizer gives exactly the full \(S^\Omega\)-image.
Together with the onto base map this is surjectivity
to \(E\). Conversely every onto map has precisely
these coordinate images and distinct kernels.
The gauge factor is the literal one from
\textup{(D.47.17)}. The given epimorphism \(q_E\)
has these properties, so \(e_0>0\). This proves the
complete fibre statement.

For completeness the ordinary conjugator union also has
an exact comparison. Let \(\mathcal C_i\) be the
\(A\)-class of \(X_i\), and let \(r_i\) be the size of
its orbit under \(C_D(d_i)\). This action is defined by
conjugating with any lift of the centralizer element;
two such lifts differ by \(A\), so the action on
\(A\)-classes is well-defined. Conjugation by \(H_i\)
identifies that orbit with the \(A\)-classes of all
\(E\)-conjugates of \(X_i\) above \(q_1(x_i)\).
Every class is obtained by replacing the specified base
conjugator \(h_i\) by \(h_ic_i\), with
\(c_i\in C_D(d_i)\). Different triples of these
\(A\)-classes give disjoint literal source fibres.
The established result applies to every one of these
triples and gives the same cardinality
\(|S|^{|\Omega|-1}e_0\). Consequently the complete
ordinary peripheral epimorphism fibre above either
base map has cardinality
\(r_1r_2r_3|S|^{|\Omega|-1}e_0\).
This proves the stated full-union assertion with all
base centralizer choices retained.

\paragraph{The specified pro-five and coupled sources.}
Use the actual D.37 source \(q_0=\rho_n\),
\(q_1=\beta_Q\rho_n\) and its displayed conjugators
\(1,f_Q,h_Q\). Equation \textup{(D.47.3)} supplies
the asserted generating lift above those exact maps.
Every homomorphism from \(S^\Omega\) to a finite
five-group is trivial: its image is both perfect and
a finite five-group, while every nontrivial finite
five-group has nontrivial abelianization. Hence
every such homomorphism from \(E\) factors uniquely
through \(P_n\), proving the stated maximal pro-five
quotient with its actual induced action.

More generally, let \(L\twoheadrightarrow D\) have
solvable kernel \(B\), and retain a given epimorphism
\(q_L:\widehat F\twoheadrightarrow L\) above \(q_1\).
Combine the constructed \(f\) and \(q_L\) in
\(E\times_D L\). The \(K_1\)-image projects onto
both \(S^\Omega\) and \(B\). Its Goursat common
quotient is a solvable quotient of the perfect group
\(S^\Omega\), hence is trivial. Thus that image is
\(S^\Omega\times B\). The base image is \(D\),
so the combined source is onto the actual fibre
product and retains \(q_L\) literally.

\paragraph{Exact composite and nonsplit checks.}
The program \texttt{a\_cyclic\_outer\_composite.py}
retains all forty literal corrected push maps and
their inverses used in the twelve recorded moduli.
It verifies the complete nonunit-character and
boundary-supplier formulas, and all \(359268\)
genus-one handle/boundary characters in its orbit
checks. The proof above supplies the arbitrary
modulus and genus comparison.

The independent program
\texttt{a\_cyclic\_outer\_nonsplit.py}
uses the actual quotient \(Q_{16}\times C_3\to C_6\)
with kernel \(Q_8\). In coordinates its law and
projection are
\[
 \begin{split}
 (a,b,c)(x,y,z)&=
 (a+(-1)^b x+4by\bmod8,\ b+y\bmod2,\ c+z\bmod3),\\
 \epsilon(a,b,c)&=3(a\bmod2)+4c\pmod6.
 \end{split}
 \tag{D.47.25}
\]
The chosen lift \((1,0,1)\) has order \(24\);
every lift of \(1\in C_6\) has order \(12\) or \(24\),
so this quotient has no section. Over the actual
Heisenberg base \(\operatorname{Heis}(\mathbf Z/6)\),
the program checks all \(46656\) outer-cocycle
products, the six intermediate characters, the
actual gauge and embedding multiplications, \(804\)
nontrivial kernel carries, and \(6084\) literal
cycle-product transports. This verifies precisely
the nonsplit group identities used above; the
simple-kernel generation is proved by the displayed
normalizer and Goursat comparisons. Both programs,
their outputs, JSON data, and the nonsplit matrix
certificate are retained in \texttt{code/}.
\end{proof}
